\documentclass[a4paper]{report}
\usepackage[utf8]{inputenc}
\usepackage[portuguese]{babel}
\usepackage{hyperref}
\usepackage{a4wide}
\hypersetup{pdftitle={Sistema de Troca de Ficheiros},
pdfauthor={José Ferreira, João Teixeira, Miguel Solino},
colorlinks=true,
urlcolor=blue,
linkcolor=black}
\usepackage{subcaption}
\usepackage[cache=false]{minted}
\usepackage{listings}
\usepackage{booktabs}
\usepackage{multirow}
\usepackage{appendix}
\usepackage{tikz}
\usepackage{authblk}
\usepackage{verbatim}
\usetikzlibrary{positioning,automata,decorations.markings}

\begin{document}

\title{Troca de Ficheiros\\ 
\large Grupo Nº 34}
\author{José Ferreira (A83683) \and João Teixeira (A85504) \and Miguel Solino (A)}
\date{\today}

\begin{center}
    \begin{minipage}{0.75\linewidth}
        \centering
        \includegraphics[width=0.4\textwidth]{eng.jpeg}\par\vspace{1cm}
        \vspace{1.5cm}
        \href{https://www.uminho.pt/PT}
        {\color{black}{\scshape\LARGE Universidade do Minho}} \par
        \vspace{1cm}
        \href{https://www.di.uminho.pt/}
        {\color{black}{\scshape\Large Departamento de Informática}} \par
        \vspace{1.5cm}
        \maketitle
    \end{minipage}
\end{center}

\tableofcontents

\pagebreak

\chapter{Introdução}

O objetivo deste projeto é construir um sistema de partilha de ficheiros de 
musica, de forma a ser capaz de gerir os ficheiros de musica, bem como todos
os meta-dados a eles associados, aplicando uma arquitetura 
\textit{Cliente, Servidor} com recurso a \textit{Sockets TCP} e \textit{Threads},
aplicando assim os conhecimentos sobre Sistemas Distribuidos adquiridos nas 
aulas desta UC\@.
Como objetivo, é necessário implementar um sistema que permita a utilização 
concorrente de um servidor por vários clientes, assegurando a consistencia
de toda a informacao introduzida no sistema a qualquer cliente que esteja
a fazer uso da mesma.\\
Ao longo deste relatorio vamos descrever as nossas abordagens a estes problemas.

\chapter{Arquitetura e Solução do Projeto}

\section{Servidor}

No modulo Servidor sao tratadas todas as conexoes e pedidos dos clientes, bem
como a gestao da informacao relativa a logins e a musicas armazenadas de 
forma a que seja possivel um acesso concorrente entre as diversas 
\textit{Threads} lancadas para comportar varios clientes em simultaneo.
A rececao, envio e armazenamento dos ficheiros de musica e tambem mantido aqui,
tendo em vista uma boa gestao de recursos, tanto ao nivel de memoria bem como
ao nivel do \textit{CPU}.\\
De forma a conseguir um servidor capaz de lidar com varios clientes em
simultaneo, este e dividido em varias \textit{Threads}, sendo, para cada
cliente que se tente conectar, criada uma para receber e responder aos
pedidos vindos do cliente, uma para enviar notificacoes sobre novas musicas
adicionadas, uma para gerir os \textit{Downloads} requisitados, e uma ultima
para receber os ficheiros de musica enviados pelo respetivo cliente.

\section{Cliente}

No modulo Cliente e efetuada a ponte entre o utilizador e o servidor, sendo 
capaz de efetuar pedidos pedidos ao servidor e apresentar ao utilizador a
respetiva resposta aos pedidos efetuados. Esta preparado para complementar
o servidor no envio e rececao de ficheiros de musica.

\section{Comunicacao Cliente-Servidor}

\subsection{Pedidos e respostas aos pedidos}
Para efetuar a comunicacao entre o cliente e o servidor, fixamos
um formato de texto, inspirado no formato \textit{JSON}, que no caso
de pedidos ao servidor tem o formato \ref{inputformat}, onde o 
\textit{request\_type} corresponde ao tipo de pedido que o cliente quer
efetuar, por exemplo, para adicionar uma musica iria ser \textit{add\_music}.\\

\begin{figure}[H]
    \begin{verbatim}
        {type='request\_type', content=['obj1','obj2',...]}
    \end{verbatim}
    \caption{Formato dos pedidos}
    \label{inputformat}
\end{figure}

As respostas a pedidos segue o formato \ref{outputformat}, onde o 
\textit{status} pode ser \textit{ok} ou \textit{failed}, e o \textit{res}
e a resposta ao pedido efetuado, ou a razao pela qual nao foi possivel 
efetua-lo.

\begin{figure}[H]
    \begin{verbatim}
        {status='status', result='res'}
    \end{verbatim}
    \caption{Formato das repostas}
    \label{outputformat}
\end{figure}

Para comunicar com o cliente, linhas de texto com este formato sao transmitidos
atraves de uma socket criada para o efeito.

\subsection{Uploads e Downloads}

Para efetuar a transferencia de ficheiros de e para o servidor, e criada uma
\textit{ServerSocket} por cada cliente, numa porta aleatoria, e posteriormente
e enviada ao servidor a informacao para ele se conectar a essa \textit{Socket}.
Nesta \textit{socket} apenas sao enviados binarios correspondentes ao ficheiro
de musica, em partes de um tamanho maximo definido. Embora sejam possiveis
downloads e uploads concorrentes entre clientes, um cliente apenas pode efetuar
um download e um upload em simultaneo.

\subsection{Notificacao de novas musicas}

Quando e adicionada uma nova musica ao sistema, esta e adicionada a uma queue de
notificacoes que sera utilizada por cada \textit{thread} de notificacoes, 
enviando assim uma notificacao a todos os clientes que estejam ligados ao
servidor no momento.

\chapter{Conclusão}

Para concluir, conseguimos cumprir todos os requisitos propostos criando no processo um
sistema de troca de ficheiros de musica capaz de utilizar uma arquitetura 
\textit{Cliente, Servidor} que suporta vários de clientes.
Como trabalho futuro, gostaríamos de permitir a cada cliente ter mais do que um
download e upload a correr em simultaneo.

\end{document}
